\documentclass[12pt,a4paper]{article}
\usepackage[utf8]{inputenc}
\usepackage[T1]{fontenc}
\usepackage{amsmath,amssymb,amsthm}
\usepackage{graphicx}

\theoremstyle{plain}
\newtheorem{theorem}{Theorem}
\newtheorem{proposition}[theorem]{Proposition}
\newtheorem{corollary}[theorem]{Corollary}

\title{The Phase Symmetry $K(s)K(1-s)\in\mathbb{R}$:\\A Numerical Investigation}
\author{Andr\'e Gualberto}

\begin{document}
\maketitle

\begin{abstract}
We examine the condition $K(s)K(1-s)\in\mathbb{R}$ for
$K(s)=2^s-1$, which is algebraically equivalent to
$\Re(s)=1/2$ or $\Im(s)=n\pi/\ln2$ for integer $n$.
If a non-trivial zero of the Riemann zeta function were to
lie off the critical line, it could only do so on one of
these exceptional lines.  We present a numerical scan of
200 such lines ($n=1,\dots,200$) that finds no candidate
with $|\zeta(s)|<10^{-4}$.  The exclusion of these lines
remains an open problem.
\end{abstract}

\noindent{\it Key words and phrases:}\ Riemann hypothesis, phase symmetry,
zeta function, critical line.

\noindent{\it 2020 Mathematics Subject Classification:}\ Primary 11M26;
Secondary 11Y35.

\section{Introduction}
\label{sec:intro}

The observation that the number~$2$ decomposes as $2=1+1$, expressed
algebraically as
\[
\frac12 \times 2 = 1,
\]
suggests a {\it valence} function $K(s)=2^s-1$ that captures the
role of the base~$2$ in the functional equation of $\zeta(s)$.
The condition $K(s)K(1-s)\in\mathbb{R}$ is of particular interest
because it forces $\Re(s)=1/2$ unless $\Im(s)$ falls on a sparse
set of lines.

This note investigates whether any non-trivial zero of $\zeta(s)$
could lie on those exceptional lines.  We show algebraically that
the phase condition is equivalent to $\Re(s)=1/2$ or
$\Im(s)=n\pi/\ln2$, and numerically test 200 such lines.

\section{Algebraic Symmetry}
\label{sec:phase}

Define $K(s)=2^s-1$.  This function captures the {\it valence} of
the number~$2$ in the functional equation of $\zeta(s)$.

\begin{proposition}
\label{prop:symmetry}
$K(s)K(1-s)\in\mathbb{R}$ iff $\Re(s)=\tfrac12$ or
$\Im(s)=n\pi/\ln2$ for some $n\in\mathbb{Z}$.
\end{proposition}
\begin{proof}
Write $s=\sigma+it$.  Then
\[
K(s)K(1-s) = (2^s-1)(2^{1-s}-1) = 3 - 2^s - 2^{1-s}.
\]
Since $2^{1-s}=2\cdot2^{-s}$,
\[
\Im\bigl(K(s)K(1-s)\bigr) = -\Im(2^s + 2^{1-s}).
\]
With $2^s = e^{\sigma\ln2}\,e^{it\ln2}$,
\begin{align*}
\Im(2^s + 2^{1-s})
&= e^{\sigma\ln2}\sin(t\ln2)
   + e^{(1-\sigma)\ln2}\sin(-t\ln2)\\
&= \bigl(e^{\sigma} - e^{1-\sigma}\bigr)\ln2\;\sin(t\ln2).
\end{align*}
This vanishes iff $\sigma=1/2$ or $t\ln2=n\pi$.
\end{proof}

\begin{proposition}
\label{teo:equiv}
If $\rho$ is a non-trivial zero of $\zeta(s)$ for which
$K(\rho)K(1-\rho)\in\mathbb{R}$ and $\Im(\rho)\neq n\pi/\ln2$,
then $\Re(\rho)=\tfrac12$.  Whether any zero lies on an
exceptional line $t=n\pi/\ln2$ with $\sigma\neq1/2$ is an open
question.
\end{proposition}

\subsection*{Conjecture (Gualberto Phase Symmetry Conjecture)}

Based on the algebraic and numerical evidence presented, we propose:

\[
\boxed{
\text{If }\zeta(s)=0\text{ and }0<\Re(s)<1,\text{ then }K(s)K(1-s)\in\mathbb{R}.
}
\]

This conjecture is equivalent to the Riemann Hypothesis, since the
phase condition $K(s)K(1-s)\in\mathbb{R}$ forces $\Re(s)=1/2$
(except for the exceptional lines $t=n\pi/\ln2$, which are
numerically excluded but await a rigorous analytic proof).

\section{Numerical Scan of Exceptional Lines}
\label{sec:scan}

The lines $t=n\pi/\ln2$ are the only algebraic loophole in
Prop.~\ref{prop:symmetry}.  A non-trivial zero off the critical
line could, in principle, exist at $\sigma\neq1/2$ with
$t=n\pi/\ln2$.

We scanned 200 such lines for $n=1,\dots,200$ over the range
$\sigma=0.05,\dots,0.95$ with step~$0.05$ using the mpmath library
(30-digit precision).  For each point we evaluated
$|\zeta(\sigma+it)|$ and recorded the minimum on each line.

No candidate with $|\zeta(s)|<10^{-4}$ was detected.  The minimum
values remain orders of magnitude above zero across the entire
range of~$n$ tested.  This numerical evidence is consistent with
the conjecture that no zero lies on these lines, but it does not
constitute a proof.

An additional random sample of 2000 points with
$\sigma\in(0,1)$ and $t\in(0,50)$ confirmed that
$\Im(K(s)K(1-s))$ only vanishes on $\sigma=0.5$ or on the
exceptional lines (where $\zeta(s)\neq0$ within our precision).

\section{Variation with Base $q$}
\label{sec:q}

The symmetry extends to any real $q>0$, $q\neq1$ via
$K_q(s)=q^s-1$.  The condition $K_q(s)K_q(1-s)\in\mathbb{R}$
is algebraically equivalent to $\Re(s)=1/2$ or
$\Im(s)=n\pi/\ln q$, as the same calculation shows.  The critical
line $\sigma=1/2$ is therefore universal for all bases $q$.

\section{Discussion}
\label{sec:conc}

The phase condition $K(s)K(1-s)\in\mathbb{R}$ provides an
algebraic characterisation of the critical line $\Re(s)=1/2$:
it is the unique line (aside from the discrete set of exceptional
lines) on which the valence function $K(s)=2^s-1$ and its
reflection $K(1-s)$ have real product.

The numerical data suggest that no non-trivial zero of $\zeta(s)$
lies on an exceptional line, but a rigorous analytic proof of
this fact---and hence a proof that the phase condition is
equivalent to the Riemann Hypothesis---remains open.

\begin{thebibliography}{10}

\bibitem{edwards} H.~M.~Edwards,
  {\it Riemann's Zeta Function}, Dover, 1974.

\bibitem{titchmarsh} E.~C.~Titchmarsh and D.~R.~Heath-Brown,
  {\it The Theory of the Riemann Zeta Function},
  2nd ed., Oxford University Press, 1986.

\end{thebibliography}
\end{document}
